Primera prueba de corrección, en las dos máquinas. Comparamos una versión ingenua, que lee el puntero y después se apunta en el contador, con el código de los autores.

AddressSanitizer vigila el uso de memoria ya liberada. En el i5 atrapa a la versión ingenua con un use-after-free; en el Ultra 9, con una doble liberación. Es el mismo fallo de fondo, visto en otro momento. Al código de los autores no le encuentra nada, ni con 16 hilos.

ThreadSanitizer vigila los accesos simultáneos sin protección. Mismo resultado: la ingenua cae y la de los autores pasa limpia.

Que la ingenua caiga es importante: demuestra que los detectores estaban funcionando. Un cero sin esa prueba no valdría nada.
